\chapter{2015年安徽卷（理）}

\section{单选题}

\begin{problem}
设 \(i\) 是虚数单位，则复数 \(\frac{2\i}{1 - \i}\) 在复平面内所对应的点位于（\quad）

\choices
    {第一象限}
    {第二象限}
    {第三象限}
    {第四象限}
\answer{B}

\begin{solution}
\[
\frac{2\i}{1-\i}=\frac{2\i(1+\i)}{(1-\i)(1+\i)}=\frac{2\i+2\i^2}{2}=-1+\i
\]
该复数的实部为 \(-1\)，虚部为 \(1\)，对应的点在第二象限，故选 B.
\end{solution}
\end{problem}

\begin{problem}
下列函数中，既是偶函数又存在零点的是（\quad）

\choices
    {\(y = \cos x\)}
    {\(y = \sin x\)}
    {\(y = \ln x\)}
    {\(y = x^{2} + 1\)}
\answer{A}

\begin{solution}
\(y=\cos x\) 满足 \(\cos(-x)=\cos x\)，是偶函数，且当 \(x=\frac{\pi}{2}\) 时有 \(\cos x=0\)，所以存在零点. \(y=\sin x\) 是奇函数；\(y=\ln x\) 的定义域不是关于原点对称的集合，不能成为偶函数；\(y=x^2+1\) 虽是偶函数，但始终大于 \(0\)，没有零点. 因此选 A.
\end{solution}
\end{problem}

\begin{problem}
设 \(p: 1 < x < 2\)，\(q: 2^x > 1\)，则 \(p\) 是 \(q\) 成立的（\quad）

\choices
    {充分不必要条件}
    {必要不充分条件}
    {充分必要条件}
    {既不充分也不必要条件}
\answer{A}

\begin{solution}
由 \(1<x<2\) 可知 \(x>0\)，而底数 \(2>1\)，所以 \(2^x>1\)，即 \(p\Rightarrow q\). 但 \(x=3\) 时 \(2^x>1\) 成立而 \(1<x<2\) 不成立，故 \(q\) 不能推出 \(p\). 因此 \(p\) 是 \(q\) 成立的充分不必要条件，选 A.
\end{solution}
\end{problem}

\begin{problem}
下列双曲线中，焦点在 \(y\) 轴上且渐近线方程为 \(y = \pm 2x\) 的是（\quad）

\choices
    {\(x^{2} - \frac{y^{2}}{4} = 1\)}
    {\(\frac{x^2}{4} - y^2 = 1\)}
    {\(\frac{y^2}{4} - x^2 = 1\)}
    {\( y^{2}-\frac{x^{2}}{4}=1 \)}
\answer{C}

\begin{solution}
焦点在 \(y\) 轴上的双曲线应写成
\(\dfrac{y^2}{a^2}-\dfrac{x^2}{b^2}=1\)，其渐近线为 \(y=\pm\dfrac{a}{b}x\). 选项 C 中 \(a=2\)，\(b=1\)，渐近线为 \(y=\pm2x\)，且焦点在 \(y\) 轴上. 选项 D 的渐近线为 \(y=\pm\frac12x\)，其余两项焦点在 \(x\) 轴上. 故选 C.
\end{solution}
\end{problem}

\begin{problem}
已知 \(m\)，\(n\) 是两条不同直线，\(\alpha\)，\(\beta\) 是两个不同平面，则下列命题正确的是（\quad）

\choices
    {若 \(\alpha\)，\(\beta\) 垂直于同一平面，则 \(\alpha\) 与 \(\beta\) 平行}
    {若 \(m\)，\(n\) 平行于同一平面，则 \(m\) 与 \(n\) 平行}
    {若 \(\alpha\)，\(\beta\) 不平行，则在 \(\alpha\) 内不存在与 \(\beta\) 平行的直线}
    {若 \(m\)，\(n\) 不平行，则 \(m\) 与 \(n\) 不可能垂直于同一平面}
\answer{D}

\begin{solution}
选项 A 不正确. 例如取 \(\alpha:x=0\)、\(\beta:y=0\)、公共平面为 \(z=0\)，则 \(\alpha\) 与 \(\beta\) 都垂直于 \(z=0\)，但二者相交而不平行.

选项 B 不正确. 取平面 \(\pi:z=0\)，直线 \(m:y=0\)，\(z=1\) 与 \(n:x=0\)，\(z=2\) 都平行于 \(\pi\)，但两直线方向不同，故不平行.

选项 C 不正确. 例如取 \(\alpha:z=0\)、\(\beta:y=0\)，两平面相交但直线 \(y=1\)，\(z=0\) 在 \(\alpha\) 内且平行于 \(\beta\).

若两条直线都垂直于同一平面，则它们的方向都与该平面的法向量平行，从而两直线平行. 因此当 \(m\)，\(n\) 不平行时，它们不可能同时垂直于同一平面，选 D.
\end{solution}
\end{problem}

\begin{problem}
若样本数据 \(x_{1}\)，\(x_{2}\)，\(\cdots\)，\(x_{10}\) 的标准差为 8，则数据 \(2x_{1}-1\)，\(2x_{2}-1\)，\(\cdots\)，\(2x_{10}-1\) 的标准差为（\quad）

\choices
    {8}
    {15}
    {16}
    {32}
\answer{C}

\begin{solution}
设原数据的平均数为 \(\overline{x}\)，则变换后数据的平均数为 \(2\overline{x}-1\). 因而变换后每个数据与新平均数之差为
\[
(2x_i-1)-(2\overline{x}-1)=2(x_i-\overline{x})
\]
标准差是偏差平方平均数的算术平方根，所以标准差扩大为原来的 \(|2|=2\) 倍，得到 \(2\times8=16\). 故选 C.
\end{solution}
\end{problem}

\begin{problem}
一个四面体的三视图如图所示，则该四面体的表面积是（\quad）

\bitmapfigure[max width=0.47\linewidth]{img/2015/anhui_science/q7_fig1.png}

\choices
    {\(1 + \sqrt{3}\)}
    {\(2 + \sqrt{3}\)}
    {\(1 + 2\sqrt{2}\)}
    {\(2\sqrt{2}\)}
\answer{B}

\begin{solution}
根据三视图建立坐标系，可取四面体的四个顶点为
\(A=(0,0,1)\)，\(B=(-1,0,0)\)，\(C=(1,0,0)\)，\(D=(0,1,0)\). 这与正视图、侧视图和俯视图中的长度及投影关系一致.

于是
\(AB=AC=AD=BD=CD=\sqrt{2}\)，\(BC=2\).
三角形 \(ABC\) 和 \(BCD\) 的底边均为 \(2\)，高均为 \(1\)，所以它们的面积均为 \(1\). 三角形 \(ABD\) 和 \(ACD\) 都是边长为 \(\sqrt{2}\) 的等边三角形，面积均为
\(\dfrac{\sqrt{3}}{4}(\sqrt{2})^2=\dfrac{\sqrt{3}}{2}\).

因此四面体的表面积为
\[
S=1+1+\frac{\sqrt{3}}{2}+\frac{\sqrt{3}}{2}=2+\sqrt{3}
\]
故选 B.
\end{solution}
\end{problem}

\begin{problem}
\(\triangle ABC\) 是边长为 2 的等边三角形，已知向量 \(\bs{a}\)，\(\bs{b}\) 满足，\(\overrightarrow{AB} = 2\bs{a}\)，\(\overrightarrow{AC} = 2\bs{a} + \bs{b}\)，则下列结论正确的是（\quad）

\choices
    {\(|\bs{b}| = 1\)}
    {\(\bs{a} \perp \bs{b}\)}
    {\(\bs{a} \cdot \bs{b} = 1\)}
    {\((4\bs{a} + \bs{b})\perp \overrightarrow{BC}\)}
\answer{D}

\begin{solution}
由 \(\overrightarrow{BC}=\overrightarrow{AC}-\overrightarrow{AB}=\bs{b}\)，且三角形为边长 \(2\) 的等边三角形，得
\(|2\bs{a}|=2\)，\(|\bs{b}|=|\overrightarrow{BC}|=2\).
又由 \(|2\bs{a}+\bs{b}|=2\)，并且 \(|\bs{a}|=1\)，可得
\(4+4\bs{a}\cdot\bs{b}+4=4\)，\(\bs{a}\cdot\bs{b}=-1\).
所以 A、B、C 均不正确. 最后
\[
(4\bs{a}+\bs{b})\cdot\overrightarrow{BC}=(4\bs{a}+\bs{b})\cdot\bs{b}=4(-1)+4=0
\]
故 \((4\bs{a}+\bs{b})\perp\overrightarrow{BC}\)，选 D.
\end{solution}
\end{problem}

\begin{problem}
函数 \( f(x) = \frac{ax + b}{(x + c)^2} \) 的图象如图所示，则下列结论成立的是（\quad）

\bitmapfigure[max width=0.30\linewidth]{img/2015/anhui_science/q9_fig1.png}

\choices
    {\(a > 0\)，\(b > 0\)，\(c < 0\)}
    {\(a < 0\)，\(b > 0\)，\(c > 0\)}
    {\(a < 0\)，\(b > 0\)，\(c < 0\)}
    {\(a < 0\)，\(b < 0\)，\(c < 0\)}
\answer{C}

\begin{solution}
由图象的竖直渐近线在 \(y\) 轴右侧，且渐近线为 \(x=-c\)，可知 \(-c>0\)，所以 \(c<0\). 图象在 \(x=0\) 处位于 \(x\) 轴上方，而
\(f(0)=\dfrac{b}{c^2}>0\)，故 \(b>0\). 图象与 \(x\) 轴的交点在正半轴上，零点 \(-\dfrac{b}{a}>0\)，结合 \(b>0\) 得 \(a<0\). 因此只有选项 C 正确.
\end{solution}
\end{problem}

\begin{problem}
已知函数 \( f(x) = A \sin (\omega x + \varphi) \) （\(A\)，\(\omega\)，\(\varphi\) 均为正的常数） 的最小正周期为 \( \pi \)，当 \( x = \frac{2\pi}{3} \) 时，函数 \( f(x) \) 取得最小值，则下列结论正确的是（\quad）

\choices
    {\(f(2) < f(-2) < f(0)\)}
    {\(f(0) < f(2) < f(-2)\)}
    {\(f(-2) < f(0) < f(2)\)}
    {\(f(2) < f(0) < f(-2)\)}
\answer{A}

\begin{solution}
最小正周期满足 \(\dfrac{2\pi}{\omega}=\pi\)，所以 \(\omega=2\). 在 \(x=\frac{2\pi}{3}\) 时取得最小值，故
\[
\sin\left(\frac{4\pi}{3}+\varphi\right)=-1
\]
从而 \(\varphi=\frac{\pi}{6}+2k\pi\). 相位相差整数个周期不改变函数，故可取
\(f(x)=A\sin\left(2x+\frac{\pi}{6}\right)\).

因为 \(7\pi/6<4<4\pi/3\)，所以 \(\sin4<0\)，并且
\(\frac{\pi}{6}-4\) 加上 \(2\pi\) 后落在 \((5\pi/6,\pi)\) 内. 于是
\[
f(2)-f(-2)=A\left[\sin\left(4+\frac{\pi}{6}\right)-\sin\left(\frac{\pi}{6}-4\right)\right]=A\sqrt{3}\sin4<0
\]
且 \(f(-2)=A\sin(\frac{\pi}{6}-4)<A\sin\frac{\pi}{6}=f(0)\). 因此
\(f(2)<f(-2)<f(0)\)，选 A.
\end{solution}
\end{problem}

\section{填空题}

\begin{problem}
\(\left(x^{3} + \frac{1}{x}\right)^{7}\) 的展开式中 \(x^{5}\) 的系数是\fillinblank{}.
\answer{35}

\begin{solution}
展开式的第 \(k+1\) 项为
\[
\mathrm C_7^k(x^3)^{7-k}\left(\frac1x\right)^k=\mathrm C_7^k x^{21-4k}
\]
令 \(21-4k=5\)，得 \(k=4\)，所以 \(x^5\) 的系数为
\(\mathrm C_7^4=35\).
\end{solution}
\end{problem}

\begin{problem}
在极坐标系中，圆 \(\rho = 8\sin \theta\) 上的点到直线 \(\theta = \frac{\pi}{3}\)（\(\rho \in \R\)）距离的最大值是\fillinblank{}.
\answer{6}

\begin{solution}
由极坐标关系 \(x=\rho\cos\theta\)、\(y=\rho\sin\theta\)，圆的方程化为
\(\rho^2=8\rho\sin\theta\)，\(\Longrightarrow\)，\(x^2+y^2=8y\)
即
\[
x^2+(y-4)^2=16
\]
所以圆心为 \(O_1=(0,4)\)，半径为 \(4\). 由于 \(\rho\in\R\)，直线 \(\theta=\frac\pi3\) 是过原点的整条直线，其方程为 \(y=\sqrt3x\). 圆心到该直线的距离为
\[
d=\frac{|\sqrt3\cdot0-4|}{\sqrt{3+1}}=2
\]
圆上点到该直线的最大距离等于圆心到直线的距离加半径，故最大值为 \(2+4=6\).
\end{solution}
\end{problem}

\begin{problem}
执行如图所示的程序框图（算法流程图），输出的 \(n\) 为\fillinblank{}.

\bitmapfigure[max width=0.52\linewidth]{img/2015/anhui_science/q13_fig1.png}
\answer{4}

\begin{solution}
初始时 \(a=1\)，\(n=1\). 第一次判断满足 \(|1-1.414|=0.414\geqslant0.005\)，更新为
\(a=1+\frac1{1+1}=1.5\)，\(n=2\).
第二次判断仍满足条件，更新为
\(a=1+\frac1{1+1.5}=1.4\)，\(n=3\).
第三次判断仍满足条件，更新为
\(a=1+\frac1{1+1.4}=1+\frac5{12}=\frac{17}{12}\)，\(n=4\).
此时
\[
\left|\frac{17}{12}-1.414\right|=\frac{17}{12}-\frac{707}{500}=\frac2{750}<0.005
\]
条件不成立，程序输出 \(n=4\).
\end{solution}
\end{problem}

\begin{problem}
已知数列 \( \{a_{n}\} \) 是递增的等比数列，\( a_{1} + a_{4} = 9 \)，\( a_{2}a_{3} = 8 \)，则数列 \( \{a_{n}\} \) 的前 \(n\) 项和等于\fillinblank{}.
\answer{\(2^n - 1\)}

\begin{solution}
设首项为 \(a_1\)，公比为 \(q\). 由 \(a_2a_3=a_1^2q^3=8>0\) 知 \(q>0\). 再由
\(a_1(1+q^3)=9\) 知 \(a_1>0\)，数列递增从而 \(q>1\).

令 \(r=q^3\)，则
\(a_1^2r=8\)，\(a_1(1+r)=9\).
消去 \(a_1\) 得
\(\frac{81r}{(1+r)^2}=8\)，\(8r^2-65r+8=0\)
所以 \(r=8\) 或 \(r=\frac18\). 因为 \(r=q^3>1\)，故 \(r=8\)，从而 \(q=2\)，再由 \(a_1(1+8)=9\) 得 \(a_1=1\).

因此前 \(n\) 项和为
\[
S_n=\frac{a_1(q^n-1)}{q-1}=2^n-1
\]
\end{solution}
\end{problem}

\begin{problem}
设 \(x^{3} + ax + b = 0\)，其中 \(a\)，\(b\) 均为实数. 下列条件中，使得该三次方程仅有一个实根的是\fillinblank{}.

\begin{circlelist}
\item \(a = -3\)，\(b = -3\)；
\item \(a = -3\)，\(b = 2\)；
\item \(a = -3\)，\(b > 2\)；
\item \(a = 0\)，\(b = 2\)；
\item \(a = 1\)，\(b = 2\).
\end{circlelist}
\answer{\(\circled{1}\circled{3}\circled{4}\circled{5}\)}

\begin{solution}
令 \(g(x)=x^3+ax+b\)，则 \(g'(x)=3x^2+a\).

当 \(a=-3\) 时，\(g'(x)=3(x^2-1)\)，极值点为 \(x=-1\)，\(1\)，且
\(g(-1)=b+2\)，\(g(1)=b-2\).
条件 \(1\) 中 \(b=-3\)，两个极值点的函数值均为负，曲线只在 \(x>1\) 上与 \(x\) 轴相交一次；条件 \(2\) 中 \(b=2\)，方程为
\(x^3-3x+2=(x-1)^2(x+2)=0\)，有两个不同的实根，故不满足；条件 \(3\) 中 \(b>2\)，两个极值点的函数值均为正，曲线只在 \(x<-1\) 上与 \(x\) 轴相交一次.

条件 \(4\) 中方程为 \(x^3+2=0\)，函数 \(x^3+2\) 严格递增，只有一个实根；条件 \(5\) 中 \(g'(x)=3x^2+1>0\)，函数严格递增，也只有一个实根. 所以正确条件为 \(1\)，\(3\)，\(4\)，\(5\).
\end{solution}
\end{problem}

\section{解答题}

\begin{problem}
在 \(\triangle ABC\) 中，\(A = \frac{3\pi}{4}\)，\(AB = 6\)，\(AC = 3\sqrt{2}\)，点 \(D\) 在 \(BC\) 边上，\(AD = BD\)，求 \(AD\) 的长.
\begin{solution}
由余弦定理，
\[
BC^2=AB^2+AC^2-2AB\cdot AC\cos A=36+18-36\sqrt2\left(-\frac{\sqrt2}{2}\right)=90
\]
所以 \(BC=3\sqrt{10}\). 设 \(AD=BD=x\)，则 \(DC=3\sqrt{10}-x\). 对三角形 \(ABC\) 使用斯图尔特定理，得
\[
AC^2\cdot BD+AB^2\cdot DC=BC\left(AD^2+BD\cdot DC\right)
\]
代入已知量，
\[
18x+36(3\sqrt{10}-x)=3\sqrt{10}\left[x^2+x(3\sqrt{10}-x)\right]=90x
\]
整理得 \(108\sqrt{10}=108x\)，所以
\[
AD=x=\sqrt{10}
\]
\end{solution}
\end{problem}

\begin{problem}
已知 \(2\) 件次品和 \(3\) 件正品混放在一起，现需要通过检测将其区分，每次随机检测一件产品，检测后不放回，直到检测出 \(2\) 件次品或者检测出 \(3\) 件正品时检测结束.

\begin{enumerate}
\item 求第一次检测出的是次品且第二次检测出的是正品的概率；
\item 已知每检测一件产品需要费用 \(100\) 元，设 \(X\) 表示直到检测出 \(2\) 件次品或者检测出 \(3\) 件正品时所需要的检测费用（单位：元），求 \(X\) 的分布列和均值（数学期望）.
\end{enumerate}
\begin{solution}
\begin{enumerate}
\item 第一次检测出次品、第二次检测出正品的概率为
\[
P=\frac25\cdot\frac34=\frac3{10}
\]

\item 若前两次都是次品，第二次检测结束，所以
\[
P(X=200)=\frac25\cdot\frac14=\frac1{10}
\]
第三次结束的情形为 \(BGB\)、\(GBB\) 或 \(GGG\)，其中 \(B\) 表示次品，\(G\) 表示正品. 这三个有序结果的概率均为 \(\frac1{10}\)，因此
\[
P(X=300)=\frac3{10}
\]
前四次才结束时，前三次必须恰有一件次品和两件正品，其概率为
\[
P(X=400)=\frac{\mathrm C_2^1\mathrm C_3^2}{\mathrm C_5^3}=\frac6{10}
\]
此时第四件无论是次品还是正品，都会使相应数量达到终止条件. 因而 \(X\) 的分布为
\(P(X=200)=\frac1{10}\)、\(P(X=300)=\frac3{10}\)、\(P(X=400)=\frac6{10}\)，且
\[
E(X)=200\cdot\frac1{10}+300\cdot\frac3{10}+400\cdot\frac6{10}=350
\]
\end{enumerate}
\end{solution}
\end{problem}

\begin{problem}
设 \(n \in \N^*\)，\(x_n\) 是曲线 \(y = x^{2n+2} + 1\) 在点 \((1,2)\) 处的切线与 \(x\) 轴交点的横坐标.

\begin{enumerate}
\item 求数列 \( \{x_{n}\} \) 的通项公式；
\item 记 \(T_{n} = x_{1}^{2}x_{3}^{2}\cdots x_{2n - 1}^{2}\). 证明：\(T_{n} \geqslant \frac{1}{4n}\).
\end{enumerate}
\begin{solution}
\begin{enumerate}
\item 函数 \(y=x^{2n+2}+1\) 在 \((1,2)\) 处的导数为 \(2n+2\)，所以切线方程为
\[
y-2=(2n+2)(x-1)
\]
令 \(y=0\)，得
\(-2=(2n+2)(x_n-1)\)，\(x_n=1-\frac1{n+1}=\frac n{n+1}\).

\item 由通项公式，
\[
T_n=\prod_{k=1}^n\left(\frac{2k-1}{2k}\right)^2
\]
当 \(k\geqslant2\) 时，
\[
(2k-1)^2-2k(2k-2)=1>0
\]
故
\[
\left(\frac{2k-1}{2k}\right)^2>\frac{2k-2}{2k}=\frac{k-1}{k}
\]
当 \(n=1\) 时，\(T_1=\frac14=\frac1{4n}\)；当 \(n\geqslant2\) 时，
\[
T_n=\frac14\prod_{k=2}^n\left(\frac{2k-1}{2k}\right)^2
>\frac14\prod_{k=2}^n\frac{k-1}{k}=\frac1{4n}
\]
所以对任意 \(n\in\N^*\)，都有 \(T_n\geqslant\frac1{4n}\).
\end{enumerate}
\end{solution}
\end{problem}

\begin{problem}
如图所示，在多面体 \(A_{1}B_{1}D_{1}DCBA\)，四边形 \(AA_{1}B_{1}B\)，\(ADD_{1}A_{1}\)，\(ABCD\) 均为正方形，\(E\) 为 \(B_{1}D_{1}\) 的中点，过 \(A_{1}\)，\(D\)，\(E\) 的平面交 \(CD_{1}\) 于 \(F\).

\begin{enumerate}
\item 证明：\(EF \parallel B_{1}C\)；
\item 求二面角 \(E - A_{1}D - B_{1}\) 的余弦值.
\end{enumerate}

\bitmapfigure[max width=0.56\linewidth]{img/2015/anhui_science/q19_fig1.png}
\begin{solution}
\begin{enumerate}
\item
设三个正方形的边长为 \(1\)，建立直角坐标系：
\(A=(0,0,0)\)，\(\ B=(0,1,0)\)，\(\ C=(1,1,0)\)，\(\ D=(1,0,0)\)
\(A_1=(0,0,1)\)，\(\ B_1=(0,1,1)\)，\(\ D_1=(1,0,1)\).
因为 \(E\) 是 \(B_1D_1\) 的中点，所以 \(E=(\frac12,\frac12,1)\). 平面 \(A_1DE\) 的方程为
\[
x-y+z=1
\]
设 \(F\in CD_1\)，则可写成 \(F=(1,1-t,t)\). 代入平面方程得 \(2t=1\)，故
\(F=(1,\frac12,\frac12)\). 于是
\(\overrightarrow{EF}=\left(\frac12,0,-\frac12\right)=\frac12(1,0,-1)\)，\(\overrightarrow{B_1C}=(1,0,-1)\)
所以 \(EF\parallel B_1C\).

\item
平面 \(A_1DE\) 的一个法向量为 \(\bs{n}_1=(1,-1,1)\)，平面 \(B_1A_1D\) 的一个法向量为 \(\bs{n}_2=(1,0,1)\). 设二面角的平面角为 \(\theta\)，则
\[
\cos\theta=\frac{|\bs{n}_1\cdot\bs{n}_2|}{|\bs{n}_1||\bs{n}_2|}
=\frac{2}{\sqrt3\sqrt2}=\frac{\sqrt6}{3}
\]
\end{enumerate}
\end{solution}
\end{problem}

\begin{problem}
设椭圆 \(E\) 的方程为 \(\frac{x^2}{a^2} + \frac{y^2}{b^2} = 1\)（\(a > b > 0\)），点 \(O\) 为坐标原点. 点 \(A\) 的坐标为 \((a, 0)\)，点 \(B\) 的坐标为 \((0, b)\)，点 \(M\) 在线段 \(AB\) 上，满足 \(|BM| = 2|MA|\). 直线 \(OM\) 的斜率为 \(\frac{\sqrt{5}}{10}\).

\begin{enumerate}
\item 求 \(E\) 的离心率 \(e\)；
\item 设点 \(C\) 的坐标为 \((0, -b)\)，\(N\) 为线段 \(AC\) 的中点，点 \(N\) 关于直线 \(AB\) 的对称点的纵坐标为 \(\frac{7}{2}\)，求 \(E\) 的方程.
\end{enumerate}
\begin{solution}
因为 \(|BM|=2|MA|\)，所以 \(M\) 将 \(AB\) 按 \(2:1\) 分点，从而
\[
M=\left(\frac{2a}{3},\frac b3\right)
\]
由直线 \(OM\) 的斜率，
\(\frac{b/3}{2a/3}=\frac b{2a}=\frac{\sqrt5}{10}\)，\(a^2=5b^2\).

\begin{enumerate}
\item 椭圆的离心率为
\[
e=\frac{\sqrt{a^2-b^2}}a=\frac{\sqrt{4b^2}}a=\frac{2b}{a}=\frac{2\sqrt5}{5}
\]

\item 点 \(N\) 为 \(AC\) 的中点，所以 \(N=(\frac a2,-\frac b2)\). 直线 \(AB\) 的方程为
\[
bx+ay-ab=0
\]
设 \(N'\) 是 \(N\) 关于 \(AB\) 的对称点. 对直线 \(Ax+By+C=0\)，点 \((x_0,y_0)\) 的对称点纵坐标为
\(y_0-\frac{2B(Ax_0+By_0+C)}{A^2+B^2}\). 在本题中，代入 \(A=b\)，\(B=a\)，\(C=-ab\) 及 \(N\)，有
\[
bx_N+ay_N-ab=-ab
\]
因此
\[
y_{N'}=-\frac b2+\frac{2a^2b}{a^2+b^2}
\]
利用 \(a^2=5b^2\)，得
\[
y_{N'}=-\frac b2+\frac{10b^3}{6b^2}=\frac{7b}{6}
\]
由题设 \(\frac{7b}{6}=\frac72\)，所以 \(b=3\)，\(a=3\sqrt5\). 故椭圆方程为
\[
\frac{x^2}{45}+\frac{y^2}{9}=1
\]
\end{enumerate}
\end{solution}
\end{problem}

\begin{problem}
设函数 \( f(x) = x^{2} - ax + b \).

\begin{enumerate}
\item 讨论函数 \( f(\sin x) \) 在 \( \left(-\frac{\pi}{2}, \frac{\pi}{2}\right) \) 内的单调性并判断有无极值. 有极值时求出极值；
\item 记 \( f_{0}(x) = x^{2} - a_{0}x + b_{0} \)，求函数 \( |f(\sin x) - f_{0}(\sin x)| \) 在 \( \left[-\frac{\pi}{2}, \frac{\pi}{2}\right] \) 上的最大值 \(D\)；
\item 在上一问中，取 \(a_0 = b_0 = 0\)，求 \(z = b - \frac{a^2}{4}\) 满足 \(D \leqslant 1\) 时的最大值.
\end{enumerate}
\begin{solution}
\begin{enumerate}
\item 设 \(h(x)=f(\sin x)=\sin^2x-a\sin x+b\). 则
\[
h'(x)=(2\sin x-a)\cos x
\]
在给定区间内 \(\cos x>0\)，所以单调性由 \(2\sin x-a\) 的符号决定.

当 \(a\leqslant-2\) 时，对任意 \(x\in(-\frac\pi2,\frac\pi2)\)，有 \(2\sin x-a>0\)，故 \(h(x)\) 单调递增且无极值. 当 \(a\geqslant2\) 时，有 \(2\sin x-a<0\)，故 \(h(x)\) 单调递减且无极值. 当 \(-2<a<2\) 时，方程 \(2\sin x-a=0\) 在区间内有唯一解 \(x_0=\arcsin\frac a2\)，函数先减后增，在 \(x_0\) 处取得极小值
\[
h(x_0)=\left(\frac a2\right)^2-a\cdot\frac a2+b=b-\frac{a^2}{4}
\]

\item 令 \(t=\sin x\)，则 \(t\in[-1,1]\)，并且
\[
f(t)-f_0(t)=-(a-a_0)t+(b-b_0)
\]
因此
\[
D=\max_{-1\leqslant t\leqslant1}|-(a-a_0)t+(b-b_0)|
=\max\{|b-b_0-(a-a_0)|,|b-b_0+(a-a_0)|\}
\]
根据 \(\max\{|u-v|,|u+v|\}=|u|+|v|\)，得到
\[
D=|a-a_0|+|b-b_0|
\]

\item 此时 \(D=|a|+|b|\leqslant1\). 于是
\[
z=b-\frac{a^2}{4}\leqslant b\leqslant|b|\leqslant1-|a|\leqslant1
\]
取 \(a=0\)，\(b=1\) 时 \(D=1\)，且 \(z=1\)，所以 \(z\) 的最大值为 \(1\).
\end{enumerate}
\end{solution}
\end{problem}
